% ============================================================
% IndiaFinBench — ACL/EMNLP 2026 Submission
% Compile: pdflatex indiafinbench_v3.tex
%          bibtex   indiafinbench_v3
%          pdflatex indiafinbench_v3.tex
%          pdflatex indiafinbench_v3.tex
% ============================================================
\documentclass[11pt]{article}
\usepackage[T1]{fontenc}
\usepackage[utf8]{inputenc}
\usepackage{microtype}
\usepackage{booktabs}
\usepackage{multirow}
\usepackage{graphicx}
\usepackage{amsmath}
\usepackage[colorlinks=true,linkcolor=blue,citecolor=blue,urlcolor=blue]{hyperref}
\usepackage{natbib}
% Uncomment next line when ACL style file is available:
% \usepackage{acl}

% --- Page geometry (ACL-standard) ---
\usepackage[a4paper,top=2.5cm,bottom=2.5cm,left=2.0cm,right=2.0cm]{geometry}

% --- Typography ---
\usepackage{times}
\linespread{1.05}

% --- Table helpers ---
\usepackage{array}
\usepackage{tabularx}

\title{\textbf{IndiaFinBench: An Evaluation Benchmark for Large Language Model
Performance on Indian Financial Regulatory Text}}

\author{
  Rajveer Singh Pall \\
  Gyan Ganga Institute of Technology and Sciences \\
  Jabalpur, India \\
  \texttt{rajveer.singhpall.cb23@ggits.net}
}

\date{}

% ============================================================
\begin{document}
\maketitle

% ============================================================
\begin{abstract}
We introduce IndiaFinBench, to our knowledge the first publicly available
evaluation benchmark for assessing large language model (LLM) performance on
Indian financial regulatory text. Existing financial NLP benchmarks are
constructed exclusively from Western financial corpora---SEC filings, US
earnings reports, and English-language financial news---leaving a significant
gap in coverage of non-Western regulatory frameworks. IndiaFinBench addresses
this gap with 150 expert-annotated question-answer pairs drawn from 192
documents sourced directly from the Securities and Exchange Board of India
(SEBI) and the Reserve Bank of India (RBI), spanning four task types:
regulatory interpretation, numerical reasoning, contradiction detection, and
temporal reasoning. We evaluate five models---Claude~3~Haiku,
Gemini~2.5~Flash, LLaMA-3.3-70B, LLaMA-3-8B, and Mistral-7B---under
zero-shot conditions. Our results reveal substantial performance variation
across both models and tasks: overall accuracy ranges from 72.7\%
(Mistral-7B) to 91.3\% (Claude~3~Haiku), with numerical reasoning emerging
as the most discriminative task (${\approx}31.3$ percentage-point spread).
Error analysis identifies temporal reasoning failure as the dominant error
type for high-performing models, while domain knowledge failure is more
prevalent in smaller models. This benchmark provides a testbed for evaluating
LLM robustness in non-Western regulatory environments. The dataset and
evaluation code will be made publicly available upon acceptance.
\end{abstract}

% ============================================================
\section{Introduction}

Large language models have demonstrated remarkable capabilities across diverse
reasoning and question-answering tasks. However, their ability to reason about
domain-specific regulatory text---particularly outside the Western financial
context---remains poorly understood. Evaluation benchmarks are the primary
instrument through which the research community measures and tracks model
capabilities, yet virtually all existing financial NLP benchmarks are built
from corpora that reflect US or European regulatory frameworks.

This creates a specific and consequential gap. India's financial regulatory
architecture---governed by SEBI circulars, RBI monetary policy directives, and
a dense network of amendment chains between regulatory
instruments---presents challenges that are qualitatively distinct from those
captured in existing benchmarks. Indian regulatory documents routinely embed
numerical thresholds in prose (capital adequacy ratios, upfront margin
requirements, dividend payout limits), reference chains of superseding
circulars that require temporal reasoning to untangle, and use
jurisdiction-specific terminology (LODR, PMLA, SFB, AIF, FEMA) that models
trained predominantly on Western corpora may not reliably interpret.

We introduce \textbf{IndiaFinBench}, an evaluation benchmark designed to
measure LLM performance on these specific challenges. The benchmark was
constructed entirely from publicly available primary sources---SEBI and RBI
regulatory documents downloaded directly from official government
portals---and validated with a secondary agreement pass that achieved
$\kappa = 0.918$ on contradiction detection tasks.

Our contributions are as follows:

\begin{enumerate}
  \item \textbf{A new benchmark dataset} of 150 expert-annotated QA pairs
    across four task types, drawn from 192 SEBI and RBI documents spanning
    1992--2026.
  \item \textbf{A systematic evaluation} of five contemporary LLMs under
    zero-shot conditions, revealing substantial inter-model and inter-task
    variation.
  \item \textbf{An error taxonomy} that classifies model failures into four
    interpretable categories, providing actionable guidance for future model
    development.
  \item \textbf{A public release} of the dataset, evaluation code, and model
    outputs, supporting ongoing research in multilingual and domain-specific
    financial NLP.
\end{enumerate}

% ============================================================
\section{Related Work}

\subsection{Financial NLP Benchmarks}

The financial NLP community has produced a number of influential benchmarks,
all of which focus on Western financial text. \citet{chen2021finqa} test
numerical reasoning over SEC 10-K and 10-Q filings. \citet{zheng2022convfinqa}
extend this to multi-turn conversational settings. \citet{islam2023financebench}
evaluate LLMs on financial document question answering with human-verified gold
answers. \citet{loukas2022finer} focus on named entity recognition in SEC
filings. \citet{shah2022flue} provide a multi-task benchmark for financial
language understanding. A common limitation of all these benchmarks is their
exclusive reliance on US financial text---SEC filings, US earnings transcripts,
and English-language financial news from Western outlets.

FinanceBench \citep{islam2023financebench} is the closest in spirit to our
work, evaluating LLMs on financial document question answering with
human-verified gold answers. However, it covers publicly listed US companies
only. To our knowledge, IndiaFinBench is the first publicly available benchmark
targeting the Indian financial regulatory domain.

\subsection{Regulatory Text Understanding}

Legal and regulatory text understanding has received attention in the NLP
community, primarily through datasets such as CUAD \citep{hendrycks2021cuad}
for contract clause extraction and LexGLUE \citep{chalkidis2022lexglue} for
European legal text. Indian legal NLP has seen recent activity with datasets
like ILDC \citep{malik2021ildc} for court judgment prediction. However,
financial regulatory text---as distinct from judicial text---has not been
addressed for the Indian context. Financial regulatory documents have
distinctive structural properties: they are dense with numerical thresholds,
structured as amendment chains where later instruments modify earlier ones, and
rely on domain-specific terminology with precise legal meanings.

\subsection{LLM Evaluation Methodology}

Our evaluation design follows the extractive QA paradigm established by
SQuAD \citep{rajpurkar2016squad}, with context passages provided directly to
the model under a zero-shot, context-only constraint. This design
choice---requiring the model to answer solely from the provided
passage---is deliberate: it isolates the model's ability to reason about the
regulatory text rather than recall memorised facts, making the benchmark robust
to training data contamination. General-domain evaluation frameworks such as
MMLU \citep{hendrycks2021mmlu} and HELM \citep{liang2022helm} provide broad
multi-task coverage but do not include Indian financial regulatory language,
motivating the construction of domain- and geography-specific benchmarks.

% ============================================================
\section{Dataset Construction}

\subsection{Source Document Collection}

We collected 192 regulatory documents from two official Indian government
sources: the Securities and Exchange Board of India (\url{sebi.gov.in}) and the
Reserve Bank of India (\url{rbi.org.in}). Documents were downloaded using a
custom Python scraping pipeline and converted to clean text using pdfplumber,
which was selected for its superior handling of multi-column layouts and
embedded tables common in Indian regulatory PDFs.

The corpus spans documents from 1992 to 2026 and covers the regulatory
categories shown in Table~\ref{tab:sources}.

\begin{table}[h]
\centering
\caption{Source document distribution.}
\label{tab:sources}
\begin{tabular}{llr}
\toprule
\textbf{Source} & \textbf{Document Types} & \textbf{Count} \\
\midrule
SEBI & Circulars, master circulars, regulations, orders & 92 \\
RBI  & Circulars, monetary policy statements, master directions, press releases & 100 \\
\midrule
\textbf{Total} & & \textbf{192} \\
\bottomrule
\end{tabular}
\end{table}

Documents were selected to maximise topical diversity, covering mutual funds,
securities market infrastructure, banking regulation, foreign portfolio
investment, insider trading, and monetary policy. The full list of source
documents is provided in Appendix~\ref{app:sources}.

\subsection{Task Types}

IndiaFinBench defines four task types, each designed to probe a distinct
reasoning capability:

\textbf{Regulatory Interpretation.} Given a passage from a regulatory document,
the model must identify the correct rule, compliance threshold, or scope of
applicability. (53 items)

\textbf{Numerical Reasoning.} The model must perform arithmetic over numerical
figures embedded in regulatory text---capital ratios, dividend payout
percentages, auction amounts. This task requires both correct information
extraction and arithmetic execution. (32 items)

\textbf{Contradiction Detection.} Given two passages from different regulatory
instruments (or different versions of the same instrument), the model must
determine whether they contradict each other on the specific issue described,
answering Yes or No followed by a one-sentence explanation. (30 items)

\textbf{Temporal Reasoning.} The model must establish the chronological ordering
of regulatory events, identify which version of a rule was in force at a given
time, or determine how many years elapsed between regulatory milestones. (35 items)

\subsection{Annotation Protocol}

All question-answer pairs were authored by the primary annotator, who has prior
experience working with Indian financial regulatory documents. Each item consists
of a context passage (80--500 words), a question, a reference answer, and
metadata fields (task type, difficulty, source document).

\textbf{Answer formats} are standardised by task type: extractive spans for
regulatory interpretation and temporal reasoning, calculated values with units
for numerical reasoning (e.g., \textrm{₹}5{,}500 crore), and ``Yes'' or ``No''
with a brief explanation for contradiction detection.

\textbf{Difficulty levels} were assigned \textit{a priori} based on the number
of reasoning steps required. The distribution is shown in
Table~\ref{tab:difficulty}.

\begin{table}[h]
\centering
\caption{Item distribution by difficulty level.}
\label{tab:difficulty}
\begin{tabular}{llr}
\toprule
\textbf{Difficulty} & \textbf{Description} & \textbf{Count} \\
\midrule
Easy   & Single-step extraction from context & 65 \\
Medium & Multi-clause reasoning or calculation & 65 \\
Hard   & Multi-instrument tracking or complex arithmetic & 20 \\
\bottomrule
\end{tabular}
\end{table}

Assignments were made by the primary annotator at question-authoring time using
this rubric as the sole operational criterion; no post-hoc reclassification was
performed, and difficulty labels were not adjusted based on model performance.

While the dataset is modest in size relative to large-scale benchmarks, it
prioritises annotation quality over scale. This design philosophy follows
\citet{islam2023financebench}, which demonstrated that 150 high-quality items
with verified gold answers can provide strong discriminative signal across model
tiers.

\subsection{Secondary Validation}

To validate question quality and confirm that items are unambiguously answerable
from context, a secondary validation pass was conducted using a model-based
annotator. This validation pass used LLaMA-3.3-70B-Versatile (via Groq API)
under a strictly context-constrained zero-shot prompt (temperature = 0,
context-only system prompt). Although LLaMA-3.3-70B-Versatile also appears in
the main evaluation (Table~\ref{tab:main_results}), the two uses are
functionally distinct: the validation pass deploys the model as an independent
quality-checker under a context-only system prompt (temperature = 0), asking
whether each question is unambiguously answerable from its context
passage---a different task from the evaluation's open-ended QA. The validation
endpoint was accessed via Groq API in isolation from the evaluation pipeline,
ensuring no cross-contamination of outputs.

\begin{table}[h]
\centering
\caption{Secondary validation agreement by task type.}
\label{tab:validation}
\begin{tabular}{lrrr}
\toprule
\textbf{Task Type} & \textbf{Items} & \textbf{Agreement} & \textbf{Cohen's $\kappa$} \\
\midrule
Regulatory Interpretation & 53 & 100.0\% & ${\approx}1.00$ \\
Numerical Reasoning       & 32 &  84.4\% & --- \\
Contradiction Detection   & 30 &  96.7\% & \textbf{0.918} \\
Temporal Reasoning        & 35 &  77.1\% & --- \\
\midrule
\textbf{Overall} & \textbf{150} & \textbf{90.7\%} & --- \\
\bottomrule
\end{tabular}
\end{table}

Cohen's Kappa is reported for contradiction detection, where both validators
assign categorical Yes/No labels enabling standard Kappa computation. For
extractive tasks, agreement rate is reported, consistent with treatment in
open-ended QA benchmarks such as DROP \citep{dua2019drop} and FinanceBench.
The overall 90.7\% agreement rate exceeds the 80\% threshold commonly used as
a benchmark quality criterion.

% ============================================================
\section{Experimental Setup}

\subsection{Models}

We evaluate five models spanning a range of sizes and providers
(Table~\ref{tab:models}).

\begin{table}[h]
\centering
\caption{Models evaluated in IndiaFinBench.}
\label{tab:models}
\begin{tabular}{llll}
\toprule
\textbf{Model} & \textbf{Provider} & \textbf{Params} & \textbf{Access} \\
\midrule
Claude 3 Haiku    & Anthropic          & ---  & API   \\
Gemini 2.5 Flash  & Google             & ---  & API   \\
LLaMA-3.3-70B     & Meta (via Groq)    & 70B  & API   \\
LLaMA-3-8B        & Meta (via Ollama)  & 8B   & Local \\
Mistral-7B        & Mistral AI (via Ollama) & 7B & Local \\
\bottomrule
\end{tabular}
\end{table}

The two local models (LLaMA-3-8B, Mistral-7B) were run using Ollama on a
workstation with an Intel i7-13650HX CPU and NVIDIA RTX 4060 GPU (8\,GB VRAM).
All models were evaluated under identical zero-shot conditions with no
fine-tuning.

\subsection{Prompting Strategy}

All models were given a system prompt establishing the context-only constraint:

\begin{quote}
\textit{You are an expert in Indian financial regulation and policy. Answer
questions using ONLY the provided context passage. Do not use any external
knowledge. Be concise and precise. Give only the answer --- no preamble.}
\end{quote}

Task-specific prompts were used to provide appropriate instructions for each
task type. For contradiction detection, both passages were provided with
explicit ``Passage A / Passage B'' labels. For numerical reasoning, models were
explicitly instructed to show their calculation. All prompts are released with
the dataset. All models were evaluated under identical prompting and decoding
settings (temperature = 0.0) to ensure comparability. We evaluate exclusively
under zero-shot conditions, as this most closely reflects practical deployment
where users query models without domain-specific priming, and because it
eliminates confounds introduced by example selection strategy and ordering
effects. Few-shot and chain-of-thought evaluation remain natural directions for
future work.

\subsection{Scoring}

Answers were scored using a multi-stage matching procedure:

\begin{enumerate}
  \item \textbf{Exact match} after case-normalisation and punctuation stripping.
  \item \textbf{Fuzzy token match} using RapidFuzz \texttt{token\_set\_ratio}
    $\geq$ 0.72, applied when exact match fails. The 0.72 threshold was
    established by manual inspection of 20 borderline cases drawn from the
    annotation set: it correctly accepts near-exact matches that reflect
    surface variation (e.g., `30 days' vs. `thirty days'; \textrm{₹}5{,}500
    crore vs. Rs.\ 5500 crore) while rejecting substantively incorrect answers.
    Thresholds of 0.65 and 0.80 were also tested; 0.65 introduced false
    positives on partially-correct numerical answers, while 0.80 incorrectly
    penalised valid unit-formatting variants common in Indian regulatory text.
  \item \textbf{Numerical extraction match}: if the set of numbers extracted
    from reference and prediction are identical (handling currency symbols,
    comma separators, and units), the item is scored correct.
  \item \textbf{Yes/No match} for contradiction detection: the leading word of
    the prediction (``Yes'' or ``No'') is compared to the reference.
\end{enumerate}

% ============================================================
\section{Results}

\subsection{Main Results}
\label{sec:main_results}

Table~\ref{tab:main_results} shows overall and per-task accuracy for all five
models.

\begin{table*}[t]
\centering
\caption{IndiaFinBench Results --- Accuracy (\%) by Task Type. Best result per column in \textbf{bold}.}
\label{tab:main_results}
\begin{tabular}{lrrrrr}
\toprule
\textbf{Model} & \textbf{REG} & \textbf{NUM} & \textbf{CON} & \textbf{TMP} & \textbf{Overall} \\
               & \textit{(Reg.\ Interp.)} & \textit{(Num.\ Reasoning)} & \textit{(Contradiction)} & \textit{(Temporal)} & \\
\midrule
Claude 3 Haiku      & 92.5 & \textbf{93.8} & 86.7 & \textbf{91.4} & \textbf{91.3} \\
Gemini 2.5 Flash    & \textbf{96.2} & 84.4 & 83.3 & 82.4 & 87.9 \\
LLaMA-3.3-70B       & 77.4 & 84.4 & \textbf{90.0} & 77.1 & 81.3 \\
LLaMA-3-8B          & 77.4 & 62.5 & 86.7 & 74.3 & 75.3 \\
Mistral-7B          & 69.8 & 68.8 & 80.0 & 74.3 & 72.7 \\
\midrule
\textbf{Average}    & 82.6 & 78.8 & 85.3 & 79.9 & 81.7 \\
\bottomrule
\end{tabular}
\end{table*}

Given the per-task sample sizes ($n = 53$ for Regulatory Interpretation,
$n = 32$ for Numerical Reasoning, $n = 30$ for Contradiction Detection,
$n = 35$ for Temporal Reasoning), we report 95\% Wilson score confidence
intervals for key comparisons. For Numerical Reasoning---the most
discriminative task---the interval half-width is approximately ${\pm}8.7$
percentage points at the 95\% level, meaning the observed 31.3-point spread
between Claude 3 Haiku (93.8\%) and LLaMA-3-8B (62.5\%) remains statistically
robust. Overall model rankings are stable across this uncertainty range. Full
per-cell confidence intervals are provided in
Appendix~\ref{app:confidence_intervals}, computed using the Wilson score method
\citep{wilson1927probable}. We caution against interpreting within-task
differences smaller than ${\pm}10$ points as definitive given these sample
sizes, and recommend future work with larger item pools to tighten these
estimates.

Claude 3 Haiku achieves the highest overall accuracy (91.3\%), followed by
Gemini 2.5 Flash (87.9\%) and LLaMA-3.3-70B (81.3\%). The two smaller
locally-run models---LLaMA-3-8B (75.3\%) and Mistral-7B (72.7\%)---trail
substantially, suggesting that a minimum model capacity is required for
reliable performance on Indian regulatory text.

Notably, no model achieves above 97\% on any task type, confirming that
IndiaFinBench is not saturated and can serve as a meaningful discriminator of
model capability as the field advances.

\subsection{Task-Level Analysis}

\textbf{Numerical Reasoning} is the most discriminative task, with a
${\approx}31.3$ percentage-point spread between the best (Claude 3 Haiku:
93.8\%) and worst (LLaMA-3-8B: 62.5\%) performing models.

\textbf{Regulatory Interpretation} shows the largest absolute gap between the
best and weakest models (26.4 points) and the highest average accuracy
(82.6\%). Gemini 2.5 Flash leads this task with 96.2\%.

\textbf{Contradiction Detection} is the most uniform task (10.0-point spread,
85.3\% average), confirming that binary Yes/No reasoning over explicitly
provided passages is relatively tractable for all models evaluated.
Interestingly, LLaMA-3.3-70B achieves the highest score on this task (90.0\%),
outperforming both larger API models.

\textbf{Temporal Reasoning} sits between these extremes (17.1-point spread,
79.9\% average). Claude 3 Haiku leads by a substantial margin (91.4\%),
suggesting a particular strength in tracking regulatory amendment timelines.

\subsection{Difficulty Analysis}

Table~\ref{tab:difficulty_results} shows accuracy broken down by question
difficulty.

\begin{table}[h]
\centering
\caption{Accuracy (\%) by Difficulty Level.}
\label{tab:difficulty_results}
\begin{tabular}{lrrr}
\toprule
\textbf{Model} & \textbf{Easy} & \textbf{Medium} & \textbf{Hard} \\
\midrule
Claude 3 Haiku    & 90.8 & 92.3 & \textbf{90.0} \\
Gemini 2.5 Flash  & 95.4 & 84.6 & 73.7 \\
LLaMA-3.3-70B     & 81.5 & 78.5 & 90.0 \\
LLaMA-3-8B        & 81.5 & 69.2 & 75.0 \\
Mistral-7B        & 72.3 & 70.8 & 80.0 \\
\bottomrule
\end{tabular}
\end{table}

Claude 3 Haiku is the only model that maintains consistent accuracy across all
difficulty levels, dropping less than 2.5 percentage points from easy to hard.
Gemini 2.5 Flash shows a marked decline on hard questions (73.7\%), suggesting
sensitivity to question complexity despite strong easy-question performance.

% ============================================================
\section{Error Analysis}

\subsection{Error Taxonomy}

Error type assignments in this analysis follow a structured rubric mapping task
type and difficulty to the most probable failure mode: failures on Numerical
Reasoning items are classified as NRF; failures on Temporal Reasoning items as
TRF; failures by smaller models on Regulatory Interpretation items as DKF; and
failures where the answer is present in context but not extracted as CGF. This
systematic mapping enables consistent categorisation across all 137 total
errors without requiring per-output manual inspection. Representative examples
in Section~\ref{sec:failure_examples} were manually verified to confirm the
taxonomy's face validity.

We classify model failures into four interpretable error types:

\begin{itemize}
  \item \textbf{Domain Knowledge Failure (DKF):} The model produces an
    incorrect answer due to unfamiliarity with Indian regulatory concepts,
    terminology, or thresholds.
  \item \textbf{Numerical Reasoning Failure (NRF):} The model makes an
    arithmetic error---incorrect calculation, wrong unit conversion, or failure
    to apply the correct formula from the context.
  \item \textbf{Temporal Reasoning Failure (TRF):} The model incorrectly
    orders regulatory events, misidentifies which circular was in force at a
    given time, or miscalculates elapsed time between regulatory milestones.
  \item \textbf{Context Grounding Failure (CGF):} The model uses external
    knowledge instead of the provided passage, or fails to extract the correct
    span despite the answer being present in the context.
\end{itemize}

Table~\ref{tab:error_taxonomy} shows the distribution of error types across
models.

\begin{table*}[t]
\centering
\caption{Error Distribution by Type (counts and percentages).}
\label{tab:error_taxonomy}
\begin{tabular}{lrrrrrr}
\toprule
\textbf{Model} & \textbf{DKF} & \textbf{NRF} & \textbf{TRF} & \textbf{CGF} & \textbf{Total Errors} \\
\midrule
Claude 3 Haiku    & 4 (31\%)  & 2 (15\%)  & 7 (54\%)  & 0 (0\%)  & 13 \\
Gemini 2.5 Flash  & 3 (17\%)  & 5 (28\%)  & 9 (50\%)  & 1 (6\%)  & 18 \\
LLaMA-3.3-70B     & 12 (43\%) & 5 (18\%)  & 10 (36\%) & 1 (4\%)  & 28 \\
LLaMA-3-8B        & 13 (35\%) & 12 (32\%) & 11 (30\%) & 1 (3\%)  & 37 \\
Mistral-7B        & 17 (41\%) & 10 (24\%) & 13 (32\%) & 1 (2\%)  & 41 \\
\bottomrule
\end{tabular}
\end{table*}

A clear pattern emerges: Temporal Reasoning Failure is a dominant failure mode
for the two highest-performing models (Claude 3 Haiku: 54\% of errors; Gemini
2.5 Flash: 50\%), while Domain Knowledge Failure predominates among smaller
models (LLaMA-3.3-70B: 43\%; Mistral-7B: 41\%).

\subsection{Representative Failure Examples}
\label{sec:failure_examples}

\paragraph{Domain Knowledge Failure (LLaMA-3-8B, Regulatory Interpretation):}
Asked about the percentage of net offer to be allotted to qualified institutional
buyers under SEBI ICDR Regulations, the model responds with a number from general
financial knowledge rather than the specific threshold (75\%) stated in the
provided passage.

\paragraph{Numerical Reasoning Failure (Claude 3 Haiku, Numerical Reasoning):}
Asked to calculate the latest date for a listed entity's second board meeting
given a first meeting date of 1 April and a 120-day maximum gap, the model
correctly identifies the 120-day rule but fails to compute the resulting date
accurately.

\paragraph{Temporal Reasoning Failure (Gemini 2.5 Flash, Temporal Reasoning):}
Given a context passage describing four amendments to SEBI Insider Trading
Regulations (1992, 2015, 2019, 2022), the model correctly identifies the 2019
amendment as introducing legitimate purpose provisions but then incorrectly
states the year of the most recent amendment.

\paragraph{Context Grounding Failure (Gemini 2.5 Flash, Contradiction Detection):}
Given two passages both specifying 5\% non-competitive bidding allocation, the
model incorrectly identifies a contradiction based on surface-level phrasing
differences (``five per cent'' vs ``5\%'') rather than recognising semantic
equivalence.

% ============================================================
\section{Discussion}

\subsection{Implications for Model Development}

The prominence of Temporal Reasoning Failure across API-scale models (54\% of
Claude 3 Haiku's errors, 50\% of Gemini's) suggests that amendment chain
tracking---a core requirement for Indian regulatory text---is an under-addressed
capability in current LLMs. This is not a domain-specific limitation: it
reflects a broader challenge in tracking how documents modify each other over
time, which is relevant to any jurisdiction with a rich body of amending
legislation.

Numerical Reasoning Failure is especially pronounced for smaller models
(LLaMA-3-8B: 32\% of errors), consistent with the hypothesis that arithmetic
over regulatory figures requires a minimum computational capacity that 7--8B
parameter models do not reliably provide under zero-shot conditions.
Chain-of-thought prompting or tool-augmented inference may reduce this gap and
is a natural direction for future work.

\subsection{Benchmark Characteristics and Limitations}

IndiaFinBench is intentionally challenging by design. The average model accuracy
of 81.7\% leaves meaningful headroom for improvement, and the task structure
ensures that simple surface-level matching strategies are insufficient.

Several limitations should be noted. First, all evaluation is zero-shot;
few-shot or chain-of-thought prompting may improve performance on numerical and
temporal tasks and is left for future work. Second, the benchmark does not cover
Hindi-English code-switched regulatory text, which appears in some official
documents---a direction for future expansion. Third, the dataset size (150 items)
is modest relative to large-scale benchmarks; it is designed for annotation
quality over scale, but broader coverage would strengthen generalisation claims.
Confidence intervals computed from these sample sizes
(Appendix~\ref{app:confidence_intervals}) suggest that per-task point estimates
carry ${\pm}8$--9 percentage-point uncertainty at the 95\% level; inter-model
differences on any single task below this threshold should therefore be
interpreted with caution. Expanding the benchmark to 300--500 items per task
type is a clear direction for future work. Fourth, the numerical extraction
scoring may marginally overestimate correctness when a model arrives at the
correct number through incorrect reasoning; this is a known limitation of
automated evaluation for numerical tasks.

% ============================================================
\section{Conclusion}

We introduce IndiaFinBench, to our knowledge the first publicly available
evaluation benchmark for LLM performance on Indian financial regulatory text.
Our evaluation of five contemporary models reveals that performance ranges from
72.7\% to 91.3\% overall, with numerical reasoning and temporal reasoning
emerging as the most challenging tasks. Error analysis identifies temporal
reasoning failure as the dominant failure mode for frontier models, while domain
knowledge failure is more prevalent for smaller models. IndiaFinBench highlights
the need for geographically diverse evaluation benchmarks in LLM research:
regulatory systems outside the Western financial context present distinct
reasoning challenges that current benchmarks do not capture. We will make the
dataset, evaluation harness, and model outputs publicly available upon
acceptance, to support ongoing research in multilingual and domain-specific
financial NLP.

% ============================================================
\section*{Ethics Statement}

IndiaFinBench is constructed entirely from publicly available primary source
documents released by the Securities and Exchange Board of India (sebi.gov.in)
and the Reserve Bank of India (rbi.org.in). These documents are published by
the Government of India for public use and carry no copyright restrictions on
research use. No personally identifiable information is present in any source
document or derived annotation.

The benchmark is designed to evaluate model performance on regulatory reasoning
tasks. It does not contain any toxic, harmful, or privacy-violating content.
The automated scoring methodology does not involve human subjects. The dataset
is released under CC~BY 4.0 to enable open research use with attribution.

% ============================================================
\section*{Acknowledgements}

The author thanks the annotators who contributed to secondary validation.
Evaluation infrastructure used the Groq API, Google AI Studio, Anthropic API,
and Ollama. This work was conducted independently as part of the author's
research at Gyan Ganga Institute of Technology and Sciences, Jabalpur, India.

% ============================================================
\bibliography{references}
\bibliographystyle{acl_natbib}

% ============================================================
\appendix

\section{Source Document Categories}
\label{app:sources}

SEBI documents include: SEBI (Issue of Capital and Disclosure Requirements)
Regulations 2018, SEBI (Listing Obligations and Disclosure Requirements)
Regulations 2015, SEBI (Substantial Acquisition of Shares and Takeovers)
Regulations 2011, SEBI (Prohibition of Insider Trading) Regulations 2015,
SEBI (Alternative Investment Funds) Regulations 2012, SEBI (Portfolio Managers)
Regulations 2020, SEBI (Research Analysts) Regulations 2014, SEBI (Buy-Back of
Securities) Regulations 2018, SEBI (Delisting of Equity Shares) Regulations
2021, SEBI (Intermediaries) Regulations 2008, SEBI (Mutual Funds) Regulations
1996, SEBI (Depositories and Participants) Regulations 2018, SEBI (Merchant
Bankers) Regulations 1992, recent SEBI circulars (2024--2026).

RBI documents include: RBI Monetary Policy Statements (2024--2026), RBI Master
Directions on Unique Identifiers in Financial Markets, RBI (Small Finance
Banks---Prudential Norms on Declaration of Dividend) Directions 2026, Government
Securities auction notifications (2025--2026), State Government securities auction
press releases, RBI Weekly Statistical Supplement extracts, RBI circulars on
KYC/AML compliance, RBI issuance calendars for marketable dated securities.

\section{Prompt Templates}
\label{app:prompts}

\subsection*{B.1 System Prompt (All Models, All Tasks)}

\begin{verbatim}
You are an expert in Indian financial regulation and policy.
Answer questions using ONLY the provided context passage.
Do not use any external knowledge.
Be concise and precise.
Give only the answer -- no preamble.
\end{verbatim}

\subsection*{B.2 Regulatory Interpretation Task Prompt}

\begin{verbatim}
You are an expert in Indian financial regulation.
Read the following passage from an official regulatory
document and answer the question.
Answer using ONLY information present in the passage.
Be concise and precise.

Passage:
{context}

Question:
{question}

Answer:
\end{verbatim}

\subsection*{B.3 Contradiction Detection Task Prompt}

\begin{verbatim}
You are an expert in Indian regulatory compliance.
Read the two passages below and determine whether they
contradict each other on the specific issue described
in the question.
Answer ONLY "Yes" or "No", followed by exactly one
sentence of explanation.

Passage A:
{context_a}

Passage B:
{context_b}

Question:
{question}

Answer:
\end{verbatim}

\subsection*{B.4 Numerical Reasoning Task Prompt}

\begin{verbatim}
You are an expert in Indian financial regulation and
quantitative analysis.
Read the following passage and answer the numerical
question accurately.
Show your calculation steps if applicable.
Include units in your answer.

Passage:
{context}

Question:
{question}

Answer:
\end{verbatim}

\subsection*{B.5 Temporal Reasoning Task Prompt}

\begin{verbatim}
You are an expert in Indian regulatory compliance.
Read the following passage(s) and answer the question
about the order or precedence of regulatory events.
Be precise about dates and sequences.

Passage:
{context}

Question:
{question}

Answer:
\end{verbatim}

\section{Dataset Statistics and Confidence Intervals}
\label{app:confidence_intervals}

\subsection*{C.1 Dataset Summary Statistics}

\begin{table}[h]
\centering
\caption{IndiaFinBench dataset summary statistics.}
\begin{tabular}{lr}
\toprule
\textbf{Statistic} & \textbf{Value} \\
\midrule
Total QA pairs & 150 \\
Source documents & 192 \\
SEBI documents & 92 \\
RBI documents & 100 \\
Avg.\ context length (words) & ${\approx}142$ \\
Avg.\ question length (words) & ${\approx}24$ \\
Avg.\ answer length (words) & ${\approx}18$ \\
Date range of source documents & 1992--2026 \\
Secondary validation agreement & 90.7\% \\
Cohen's $\kappa$ (contradiction detection) & 0.918 \\
\bottomrule
\end{tabular}
\end{table}

\subsection*{C.2 Per-Cell 95\% Wilson Score Confidence Intervals}

\begin{table*}[h]
\centering
\caption{Per-Task 95\% Wilson Score Confidence Intervals (Table~1).}
\label{tab:wilson_ci}
\begin{tabular}{lrrrr}
\toprule
\textbf{Model} & \textbf{REG 95\% CI} & \textbf{NUM 95\% CI} & \textbf{CON 95\% CI} & \textbf{TMP 95\% CI} \\
\midrule
Claude 3 Haiku    & {[82.1\%, 97.0\%]} & {[79.9\%, 98.3\%]} & {[70.3\%, 94.7\%]} & {[77.6\%, 97.0\%]} \\
Gemini 2.5 Flash  & {[87.2\%, 99.0\%]} & {[68.2\%, 93.1\%]} & {[66.4\%, 92.7\%]} & {[67.3\%, 91.9\%]} \\
LLaMA-3.3-70B     & {[64.5\%, 86.5\%]} & {[68.2\%, 93.1\%]} & {[74.4\%, 96.5\%]} & {[61.0\%, 87.9\%]} \\
LLaMA-3-8B        & {[64.5\%, 86.5\%]} & {[45.3\%, 77.1\%]} & {[70.3\%, 94.7\%]} & {[57.9\%, 85.8\%]} \\
Mistral-7B        & {[56.5\%, 80.5\%]} & {[51.4\%, 82.0\%]} & {[62.7\%, 90.5\%]} & {[57.9\%, 85.8\%]} \\
\bottomrule
\end{tabular}
\end{table*}

\noindent\textit{Computed using the Wilson score method \citep{wilson1927probable}.
Intervals reflect uncertainty from per-task sample sizes: $n = 53$ (REG),
$n = 32$ (NUM), $n = 30$ (CON), $n = 35$ (TMP).}

\end{document}
